\section{Registro de decisiones de diseño}
\label{sec:registro-decisiones}

La normativa de la Escuela Politécnica Superior de Córdoba exige justificar las
decisiones adoptadas y evaluar los riesgos asociados a ellas
\cite{reglamento2024,guiaGII2025}. Este registro recoge cada decisión en el
momento en que se toma, en lugar de reconstruirla al final del trabajo. Cada
entrada documenta el contexto, las alternativas consideradas, la decisión, sus
consecuencias, los riesgos y su estado.

\begin{adr}{ADR-001}{Stack tecnológico principal: Python + Dash de extremo a extremo}{2026-09-28}
\adrfield{Contexto}{El TFG es un trabajo aplicado de análisis de datos de rendimiento de corredores. Debe cubrir ingesta, análisis, modelado predictivo y visualización con recursos y tiempo limitados (12 ECTS).}
\adrfield{Alternativas consideradas}{(a) Python + Dash de extremo a extremo; (b) Python para el análisis y Power BI para el dashboard; (c) JavaScript/React + Chart.js con un backend aparte; (d) Streamlit.}
\adrfield{Decisión}{Python + Dash, con Pandas, scikit-learn, statsmodels y Plotly.}
\adrfield{Consecuencias}{Un único lenguaje desde la ingesta hasta la interfaz; el esfuerzo se concentra en el análisis, que es el núcleo intelectual del TFG. Se descarta Power BI por su menor control y por sacar la lógica del repositorio. Se descarta React por duplicar lenguajes y añadir plomería.}
\adrfield{Riesgos}{Dash es menos flexible en interfaz que un framework de frontend; acoplamiento a Plotly para los gráficos.}
\adrfield{Estado}{Aceptada.}
\adrfield{Referencias}{Decisión registrada en \texttt{AGENTS.md} del repositorio.}
\end{adr}

\begin{adr}{ADR-002}{Fuente de datos: API de Strava en primer lugar, datasets simulados como respaldo}{2026-09-28}
\adrfield{Contexto}{Se necesitan datos reales de actividad (ritmo, frecuencia cardíaca, distancia, potencia, cadencia).}
\adrfield{Alternativas consideradas}{(a) API de Strava; (b) Garmin Connect; (c) datasets públicos simulados; (d) datos propios exportados manualmente.}
\adrfield{Decisión}{Strava API primero; datasets simulados como respaldo para desarrollo y pruebas; Garmin solo si llega a ser necesario.}
\adrfield{Consecuencias}{Datos reales y variados; obliga a resolver OAuth y límites de peticiones; los tests nunca tocan la red.}
\adrfield{Riesgos}{Dependencia de términos de uso y límites de la API; posible falta de acceso a datos propios suficientes; el respaldo simulado debe ser realista y documentado para no falsear resultados.}
\adrfield{Estado}{Aceptada.}
\adrfield{Referencias}{\texttt{AGENTS.md}; propuesta de TFG (Recursos previstos).}
\end{adr}

\begin{adr}{ADR-003}{Almacenamiento: SQLite en desarrollo con esquema portable a PostgreSQL}{2026-09-28}
\adrfield{Contexto}{El TFG necesita persistencia de actividades y métricas sin añadir infraestructura que consuma tiempo.}
\adrfield{Alternativas consideradas}{(a) SQLite; (b) PostgreSQL; (c) ficheros planos (CSV/Parquet); (d) base de datos en la nube.}
\adrfield{Decisión}{SQLite en desarrollo, manteniendo el esquema relacional portable a PostgreSQL.}
\adrfield{Consecuencias}{Cero configuración, un único fichero, reproducible y versionable al margen del repositorio; la portabilidad exige evitar dialectos específicos.}
\adrfield{Riesgos}{SQLite no soporta concurrencia de escritura; migrar a PostgreSQL puede revelar incompatibilidades si se usan extensiones propias.}
\adrfield{Estado}{Aceptada.}
\adrfield{Referencias}{\texttt{AGENTS.md}.}
\end{adr}

\begin{adr}{ADR-004}{Metodología de desarrollo: TDD con pytest, alcance de tests unitarios}{2026-09-28}
\adrfield{Contexto}{El TFG debe ser reproducible y sus números verificables; se descarta SDD por coste/beneficio en un trabajo individual.}
\adrfield{Alternativas consideradas}{(a) TDD (Red--Green--Refactor) con pytest; (b) SDD con artefactos de especificación; (c) desarrollo sin tests.}
\adrfield{Decisión}{TDD con pytest, ciclo Red--Green--Refactor, alcance de tests unitarios; tests de integración solo si un problema real lo exige y previo acuerdo.}
\adrfield{Consecuencias}{Las fórmulas quedan verificadas contra valores calculados a mano; el autor escribe los tests y el código, con el agente como revisor; mayor coste inicial.}
\adrfield{Riesgos}{Disciplina requerida; riesgo de tests acoplados a la implementación (deben verificar comportamiento, no implementación); los tests nunca usan red, datos reales ni la base de datos real.}
\adrfield{Estado}{Aceptada.}
\adrfield{Referencias}{\texttt{AGENTS.md}.}
\end{adr}

\begin{adr}{ADR-005}{Entorno de desarrollo aislado con \texttt{venv} (Python 3.12.3)}{2026-09-28}
\adrfield{Contexto}{Ubuntu 24.04 aplica PEP 668, por lo que \texttt{pip install} global falla por diseño.}
\adrfield{Alternativas consideradas}{(a) \texttt{venv}; (b) \texttt{conda}; (c) \texttt{pipx}; (d) \texttt{--break-system-packages}.}
\adrfield{Decisión}{Entorno virtual \texttt{.venv} en la raíz del repositorio con Python 3.12.3, aislado del sistema (\texttt{include-system-site-packages = false}).}
\adrfield{Consecuencias}{Dependencias reproducibles y aisladas; \texttt{.venv} queda fuera del control de versiones.}
\adrfield{Riesgos}{Obliga a activar el entorno antes de ejecutar; la versión de dependencias debe fijarse en \texttt{requirements.txt}.}
\adrfield{Estado}{Aceptada.}
\adrfield{Referencias}{\texttt{AGENTS.md}.}
\end{adr}

\begin{adr}{ADR-006}{Control de versiones y flujo de trabajo: Git + GitHub con commits convencionales}{2026-09-28}
\adrfield{Contexto}{El repositorio forma parte de la presencia pública del TFG y debe reflejar su estado.}
\adrfield{Alternativas consideradas}{(a) Git + GitHub público; (b) repositorio privado; (c) sin control de versiones.}
\adrfield{Decisión}{Git + GitHub, commits convencionales en inglés e imperativos, modelo de ramas \texttt{main} / \texttt{feat/<nombre>} / \texttt{fix/<nombre>} / \texttt{docs/<nombre>}, flujo rama $\rightarrow$ commits $\rightarrow$ push $\rightarrow$ pull request $\rightarrow$ merge.}
\adrfield{Consecuencias}{Trazabilidad y evidencia del trabajo; el PR actúa como revisión incluso en solitario; nunca se suben secretos ni datos crudos.}
\adrfield{Riesgos}{Riesgo de publicar datos personales o secretos si se descuida \texttt{.gitignore}; la historia compartida no se reescribe nunca.}
\adrfield{Estado}{Aceptada.}
\adrfield{Referencias}{\texttt{AGENTS.md}.}
\end{adr}

\begin{adr}{ADR-007}{Tipología del TFG y normativa aplicable}{2026-09-28}
\adrfield{Contexto}{El TFG es de tipo aplicado y debe ajustarse a la normativa de la EPSC.}
\adrfield{Alternativas consideradas}{Las cuatro tipologías del Reglamento de TFG de la EPSC.}
\adrfield{Decisión}{Tipología b) ``Análisis y resolución de casos prácticos reales en el ámbito de la ingeniería'', que se materializa en un proyecto ``llave en mano'' que incluye el desarrollo de la solución. Normativa de referencia: Reglamento 24/2024 de TFG de la EPSC, Guía de desarrollo del TFG del Grado en Ingeniería Informática 2025/26 y guía docente de la asignatura.}
\adrfield{Consecuencias}{La memoria debe seguir el índice mínimo de esa tipología (Introducción y contexto; Objetivos; Antecedentes; Análisis del problema; Diseño; Conclusiones; Bibliografía) más Manual de usuario y Manual de código.}
\adrfield{Riesgos}{La normativa no fija formato tipográfico ni extensión, por lo que esos criterios deben confirmarse con el tutor; los aspectos formales suponen el 20~\% de la calificación.}
\adrfield{Estado}{Aceptada.}
\adrfield{Referencias}{\cite{reglamento2024,guiaGII2025}.}
\end{adr}

\begin{adr}{ADR-008}{Cadena de herramientas de la memoria: LaTeX con \texttt{latexmk}, \texttt{biber} y \texttt{biblatex}}{2026-09-28}
\adrfield{Contexto}{La memoria es un entregable central que debe ser reproducible y mantenible durante todo el TFG.}
\adrfield{Alternativas consideradas}{(a) LaTeX con \texttt{latexmk} + \texttt{biber} + \texttt{biblatex}; (b) procesador de textos (Word/LibreOffice); (c) Markdown con conversión a PDF; (d) Typst.}
\adrfield{Decisión}{LaTeX con \texttt{latexmk} como driver, \texttt{biber} y \texttt{biblatex} para la bibliografía, un fichero por capítulo y todo el formato aislado en \texttt{preamble/format.tex}.}
\adrfield{Consecuencias}{Control total del formato, numeración y referencias cruzadas automáticas, contenido versionable como texto plano y trazable en Git; exige aprender LaTeX.}
\adrfield{Riesgos}{Mayor curva de aprendizaje; el estilo de cita (APA frente a IEEE) queda abierto y debe decidirse; si el tutor exige una plantilla propia, el formato debe poder cambiarse desde un único fichero.}
\adrfield{Estado}{Aceptada (con una decisión pendiente: el estilo de cita).}
\adrfield{Referencias}{\texttt{AGENTS.md}.}
\end{adr}

\pendiente{el estilo de cita (APA frente a IEEE) continúa sin decidir. Ambos
paquetes están instalados; el cambio consiste en modificar el valor
\texttt{style=} en \texttt{preamble/packages.tex}. Debe confirmarse con el tutor
y actualizar este registro (ADR-008) cuando se decida.}
